% =====================================================================
%  General introduction (write it last, or nearly last)
% =====================================================================
\unnumberedchapter{General Introduction}
\label{chap-intro}

\section*{Context}
\todonote{Digitalisation of teaching, widespread use of LMSs (including
Moodle in Algerian universities), arrival of LLMs in education.}

\section*{Problem statement}
\todonote{LMSs deliver the same path to every learner; their built-in AI
tools answer one-off requests and model neither the learner nor the
context. Central question: how can an LLM-based, context-aware multi-agent
architecture personalize learning and help teachers improve their
materials, while keeping humans in the loop?}

\section*{Objectives}
\todonote{List 3 to 4 measurable objectives.}

\section*{Contributions}
\todonote{Summarize the contributions (Chapter~3) and the implementation
(Chapter~4).}

\section*{Thesis outline}
This thesis is organized into four chapters. \Cref{chap-lms} presents
learning management systems, their interoperability standards, the main
existing platforms and their limits with respect to adaptive learning.
\Cref{chap-sota} reviews the state of the art in artificial intelligence,
large language models and multi-agent systems applied to education, then
analyzes related work. \Cref{chap-contribution} describes the proposed
multi-agent framework and its architecture. \Cref{chap-implementation}
details the design and implementation, including the integration with
Moodle. A general conclusion summarizes the results and outlines future
work.
